% ================================================================
% Chapter 12: Stable Matching and Cultural Compatibility
% ================================================================

\chapter{Stable Matching and Cultural Compatibility}
\label{ch:matching}

\epigraph{The prohibition of incest is less a rule prohibiting marriage with
the mother, sister, or daughter, than a rule obliging the mother, sister, or
daughter to be \emph{given} to others.  It is the supreme rule of the gift.}%
{Claude Lévi-Strauss, \textit{Les Structures élémentaires de la parenté} (1949)}

% ----------------------------------------------------------------
\section{Introduction}
% ----------------------------------------------------------------

Two-sided matching is ubiquitous in human social organization.  Marriage
exchange---the foundational institution of kinship---pairs lineages on the
basis of compatibility.  Student--mentor pairing in academic apprenticeship,
patron--client relationships in Mediterranean and Mesoamerican societies,
and interfaith dialogue all involve matching ``senders'' with ``receivers''
on the basis of \emph{cultural compatibility}.

In Ideatic Diffusion Theory, compatibility reduces to resonance.  Two ideas
are compatible when they evoke each other; a match is fruitful when the
composition of the paired ideas achieves significant depth.  This chapter
formalizes these notions and proves eighteen theorems that connect matching
theory to Lévi-Strauss's kinship structures, Malinowski's Kula ring,
Mauss's gift economy, and Bourdieu's marriage strategies.

The central contribution is a suite of theorems showing that compatibility is
reflexive, symmetric, and preserved under composition---precisely the
algebraic properties that make stable matching \emph{possible}.  Without
these properties, no stable assignment of senders to receivers could exist.
Lévi-Strauss's intuition that kinship is a ``total social fact'' finds its
formal expression in the closure of compatibility under the monoid
operation.\footnote{Lévi-Strauss's notion of kinship as \emph{fait social
total} (following Mauss) implies that every aspect of social life---economic,
ritual, juridical---is implicated in the matching structure.  Our algebraic
closure results formalize this totality.}

\subsection{Historical Foundations of Matching Theory}

The mathematical theory of stable matching originates with David Gale and
Lloyd Shapley's seminal 1962 paper, ``College Admissions and the Stability
of Marriage.''  Gale and Shapley proved that for any set of men and women
with complete and transitive preference orderings, a stable matching
exists---one in which no unmatched pair mutually prefers each other to their
assigned partners.  Their deferred acceptance algorithm, elegant in its
simplicity, showed that stability is not merely a theoretical possibility but
a constructively achievable outcome.

The practical significance of stable matching was recognized by the Nobel
Committee in 2012, when Lloyd Shapley (jointly with Alvin Roth) received the
Sveriges Riksbank Prize in Economic Sciences.  Roth's contribution was to
demonstrate that stable matching theory could be applied to real-world
institutional design: the National Resident Matching Program (NRMP) for
medical residencies, school choice mechanisms in New York and Boston, and
kidney exchange programs.  Roth's work showed that abstract mathematical
properties---stability, strategy-proofness, efficiency---have direct
consequences for human welfare.\footnote{Alvin E.\ Roth, ``The Economist as
Engineer: Game Theory, Experimentation, and Computation as Tools for Design
Economics,'' \textit{Econometrica} 70(4), 2002, pp.\ 1341--1378.}

The anthropological study of matching long predates its mathematical
formalization.  Claude Lévi-Strauss's \textit{Les Structures élémentaires
de la parenté} (1949) is, in essence, a theory of marriage matching.
Lévi-Strauss distinguished two fundamental types of exchange:
\emph{restricted exchange} (échange restreint), in which two groups
directly swap marriage partners, and \emph{generalized exchange} (échange
généralisé), in which three or more groups form a directed cycle of
partner transfer.  In restricted exchange, if clan A gives women to clan B,
then clan B gives women to clan A---a bilateral matching.  In generalized
exchange, A gives to B, B gives to C, C gives to A---a cyclic matching
that cannot be decomposed into bilateral pairs.

Lévi-Strauss's analysis was deeply influenced by Marcel Mauss's theory of
the gift (\textit{Essai sur le don}, 1925).  Mauss argued that exchange is
the foundation of social solidarity: the obligation to give, to receive,
and to reciprocate constitutes the ``total social fact'' (fait social total)
that binds individuals into society.  For Lévi-Strauss, marriage exchange is
the paradigmatic instance of Mauss's triple obligation: lineages are obliged
to give women (the prohibition of incest), to receive women (exogamy), and
to reciprocate (the alliance system).  Our formal notion of compatibility---
defined as resonance between sender and receiver---captures the structural
precondition for this exchange: two lineages can exchange only if their
totemic symbols resonate.

The cross-cousin marriage systems that preoccupied Lévi-Strauss provide
particularly elegant illustrations of matching theory.  In \emph{Dravidian
kinship} (exemplified by many South Indian and Australian Aboriginal
societies), the marriageable category is defined by a binary partition of
kin into ``parallel'' (same-sex sibling links) and ``cross'' (opposite-sex
sibling links) relatives.  Ego must marry a cross-cousin---specifically,
in the most restricted systems, a bilateral cross-cousin (mother's
brother's daughter, who is also father's sister's daughter).  This system
generates a stable matching between two exogamous moieties, each moiety
constituting a ``sender'' pool for the other.

The Kariera system of Western Australia, analyzed by A.\ R.\
Radcliffe-Brown in 1913, provides the simplest four-section system:
society is divided into four named sections, and marriage rules specify
exactly which section may marry which.  The Kariera system is equivalent
to a bipartite graph with two moieties and a deterministic matching rule
within each generation.  In our framework, the four sections correspond
to four ideas $a_1, a_2, a_3, a_4 \in \IS$, with compatibility
$a_i \res a_j$ holding if and only if $i$ and $j$ belong to
marriageable sections.  The reflexivity, symmetry, and composition-preservation
of compatibility (Theorems~\ref{thm:self-compat}--\ref{thm:comp-preserves})
ensure that the Kariera system is algebraically well-defined across
generations.

Modern matching markets extend these ancient patterns into new domains.
Dating applications (Tinder, Hinge, Bumble) are literal matching markets
in which individuals present signals (profiles) and express compatibility
judgments (swipes).  The underlying algorithms---while proprietary---are
descendants of the Gale-Shapley framework, adapted for incomplete
information and stochastic preferences.  In our formalism, a dating app
user's profile is a signal $s \in \IS$, their preferences constitute a
repertoire $\mathbf{R}$, and a ``match'' occurs when mutual compatibility
$s_A \res s_B$ holds.  The reflexivity of compatibility
(Theorem~\ref{thm:self-compat}) guarantees that every user is trivially
compatible with themselves---the formal reason that self-matching must be
explicitly excluded by platform design.

Organ donation matching, formalized by Roth, Sönmez, and Ünver (2004),
provides another modern instance.  Here, the ``ideas'' are biological: a
donor kidney's immunological profile is a signal, a recipient's antibody
profile is a repertoire element, and compatibility is immunological
crossmatch compatibility.  The algebraic properties we prove in this
chapter---reflexivity, symmetry, composition preservation---do \emph{not}
generally hold for immunological compatibility (crossmatch can be
asymmetric), which is precisely why kidney exchange requires more
sophisticated algorithms than simple stable matching.  Our framework thus
identifies the structural properties that distinguish culturally-mediated
matching from biologically-constrained matching.

Assortative mating---the tendency for individuals to mate with similar
partners---has been extensively documented in sociology.  Educational
homogamy (the tendency to marry partners of similar educational attainment)
has increased dramatically in the United States and Western Europe since
the 1960s, with profound consequences for income inequality.\footnote{Robert
D.\ Mare, ``Five Decades of Educational Assortative Mating,''
\textit{American Sociological Review} 56(1), 1991, pp.\ 15--32.}
In our framework, assortative mating corresponds to a preference for
matches in which sender and receiver have similar depth:
$|\depth(s) - \depth(r)| \le \epsilon$ for some threshold $\epsilon$.
The composition quality bound (Theorem~\ref{thm:quality-bound}) implies
that assortatively matched pairs achieve quality close to
$2 \cdot \depth(a)$ (when $\depth(s) \approx \depth(r) \approx \depth(a)$),
while disassortatively matched pairs (one deep, one shallow) achieve
quality bounded by $\depth(\text{deep}) + \depth(\text{shallow})$---a
larger bound but with diminishing marginal returns from the shallow
partner's contribution.

% ----------------------------------------------------------------
\section{Core Definitions}
\label{sec:ch12-defs}
% ----------------------------------------------------------------

\begin{definition}[Compatibility]
\label{def:compatible}
Two ideas $a, b \in \IS$ are \emph{compatible} when they resonate:
\[
  \operatorname{compatible}(a, b)
  \;\;\stackrel{\text{def}}{=}\;\;
  a \res b.
\]
In kinship theory, this is the analogue of \emph{marriageability}: two
lineages are compatible when their totemic symbols evoke each other.
\end{definition}

\begin{definition}[Match Pair]
\label{def:match-pair}
A \emph{match pair} is a structure $\langle \text{sender}, \text{receiver} \rangle$
pairing a sending idea (gift, signal, myth) with a receiving idea (background,
habitus, totem).  A match pair is \emph{compatible} when sender and receiver
resonate.
\end{definition}

\begin{lstlisting}[caption={MatchPair structure in Lean~4}]
structure MatchPair (I : Type*) [IdeaticSpace I] where
  sender   : I
  receiver : I

def MatchPair.isCompatible (p : MatchPair I) : Prop :=
  compatible p.sender p.receiver
\end{lstlisting}

\begin{definition}[Composition Quality]
\label{def:comp-quality}
The \emph{composition quality} of a sender--receiver pair is the depth of
their composition:
\[
  \operatorname{quality}(s, r) \;=\; \depth(s \compose r).
\]
Higher quality means deeper cultural synthesis.
\end{definition}

% ----------------------------------------------------------------
\section{Compatibility Fundamentals}
\label{sec:ch12-compat}
% ----------------------------------------------------------------

\begin{theorem}[Self-Compatibility (12.1)]
\label{thm:self-compat}
Every idea is compatible with itself: $\forall\, a \in \IS,\; a \res a$.
\end{theorem}

\begin{proof}[Proof sketch]
By reflexivity of resonance.
\end{proof}

Endogamy---marrying within one's own group---is always ``compatible'' in the
formal sense.  Lévi-Strauss argued that exogamy is culturally
\emph{mandated} precisely because endogamy is always the structural default.
The prohibition of incest is not the prohibition of something impossible but
of something \emph{too easy}: the self-match is always available, and culture
must intervene to prevent it.

\begin{remark}[Connection to Graph Theory and Bipartite Matching]
The compatibility relation $\res$, being reflexive
(Theorem~\ref{thm:self-compat}) and symmetric
(Theorem~\ref{thm:compat-symm}), defines an undirected graph on the
ideatic space $\IS$: the \emph{compatibility graph} $G_{\res} = (\IS, E)$
where $\{a, b\} \in E$ iff $a \res b$.  Reflexivity means every vertex
has a self-loop; symmetry means every edge is undirected.  A stable
matching in the sense of Gale and Shapley is then a maximum matching in
this graph, subject to preference orderings.

When the compatibility graph is bipartite---as in moiety-based marriage
systems, where society is divided into two exogamous halves---the matching
problem reduces to bipartite matching, for which efficient polynomial-time
algorithms exist (e.g., the Hopcroft-Karp algorithm).  The four-section
systems (Kariera, Aranda) correspond to bipartite graphs with additional
structure.  The eight-section systems analyzed by Lévi-Strauss correspond
to more complex graph structures whose matching properties depend on the
specific compatibility rules encoded in the kinship terminology.

This graph-theoretic perspective connects our algebraic framework to the
extensive literature on matching in combinatorial optimization, including
Hall's marriage theorem, König's theorem, and the Birkhoff-von Neumann
theorem on doubly stochastic matrices.  Each of these classical results
has an anthropological interpretation when the vertices of the graph are
cultural ideas and the edges are resonance relations.
\end{remark}

\begin{theorem}[Compatibility is Symmetric (12.2)]
\label{thm:compat-symm}
If $a$ is compatible with $b$, then $b$ is compatible with $a$:
\[
  a \res b \;\Longrightarrow\; b \res a.
\]
\end{theorem}

\begin{proof}[Proof sketch]
By symmetry of resonance.
\end{proof}

In Lévi-Strauss's \emph{restricted exchange}, if clan A can marry into clan B,
then clan B can marry into clan A.  The symmetry of marriage exchange---the
\emph{connubium}---is a structural consequence of resonance symmetry.  This
rules out asymmetric kinship systems in which marriageability is one-directional,
a constraint that Lévi-Strauss himself noted distinguishes restricted from
generalized exchange.\footnote{In \emph{generalized exchange}, asymmetry
\emph{is} allowed---clan A gives to B, B gives to C, C gives to A---but
the pairwise compatibility remains symmetric; only the \emph{flow} of
gifts is directed.}

\begin{example}[The Kula Ring: Circular Exchange and Compatibility]
Bronisław Malinowski's study of the Kula ring in the Trobriand Islands
(\textit{Argonauts of the Western Pacific}, 1922) provides a classic
illustration of matching based on compatibility.  In the Kula system,
ceremonial objects circulate among island communities in a fixed ring:
\emph{soulava} (necklaces) travel clockwise, \emph{mwali} (armbands)
travel counterclockwise.  Each exchange is between two partners---a
sender and a receiver---who are bound by a long-term Kula relationship.

In our framework, each Kula partner brings a cultural repertoire shaped
by their island's traditions, their rank within the chieftainship
hierarchy, and their personal exchange history.  A successful Kula
exchange requires compatibility: the sender's object (signal) must
resonate with the receiver's expectations (repertoire element).  A
\emph{soulava} of great prestige resonates differently in the hands of a
paramount chief ($\depth(r_{\text{chief}} \compose s_{\text{soulava}})$
is high) than in the hands of a commoner
($\depth(r_{\text{commoner}} \compose s_{\text{soulava}})$ is lower).

The symmetry of compatibility (Theorem~\ref{thm:compat-symm}) is
instantiated in the Kula by the norm of reciprocity: if chief A finds
chief B a worthy Kula partner (compatible), then chief B finds chief A
equally worthy.  The exchange relationship is necessarily bilateral, even
though the objects flow in a single direction at any given moment.
Malinowski emphasized that the Kula partnership, once established, is
``one of the permanent and overriding features of the natives'
life''---a stable match in both the colloquial and the technical sense.
\end{example}

\begin{theorem}[Void Self-Compatibility (12.3)]
\label{thm:void-self-compat}
Void is compatible with itself: $\void \res \void$.
\end{theorem}

\begin{proof}[Proof sketch]
An instance of self-compatibility (Theorem~\ref{thm:self-compat}).
\end{proof}

The null exchange---two parties exchanging nothing---is trivially compatible.
This is the structural zero-point of all reciprocity, the ``degree zero'' of
Mauss's gift economy.  It is the baseline from which all actual exchange
departs.

\begin{theorem}[Absorbed Void Preserves Compatibility (12.4)]
\label{thm:absorbed-void}
If $a \res b$, then $(a \compose \void) \res b$.
\end{theorem}

\begin{proof}[Proof sketch]
Since $a \compose \void = a$ by the void axiom, this is immediate.
\end{proof}

Appending silence to a cultural signal does not change its compatibility.
A ritual followed by silence is the same ritual---the ``afterglow'' of
ceremony is structurally inert.  The Durkheimian \emph{effervescence
collective} dissipates, but the signal's compatibility endures.

% ----------------------------------------------------------------
\section{Composition and Compatibility}
\label{sec:ch12-comp}
% ----------------------------------------------------------------

\begin{theorem}[Composition Preserves Compatibility (12.5)]
\label{thm:comp-preserves}
If $(a, a')$ and $(b, b')$ are compatible pairs, then their compositions
are also compatible:
\[
  a \res a' \;\wedge\; b \res b'
  \quad\Longrightarrow\quad
  (a \compose b) \res (a' \compose b').
\]
\end{theorem}

\begin{proof}[Proof sketch]
By the full composition-preservation of resonance
(\texttt{resonance\_compose}).
\end{proof}

If two couples are each internally compatible, their ``children'' (compositions)
are also compatible with each other.  Cultural compatibility is
\emph{heritable} through ideatic composition.  This is Lévi-Strauss's
\emph{atom of kinship}: the minimal structure of kinship---two
allied families producing compatible offspring---is formally guaranteed by
this theorem.\footnote{Lévi-Strauss introduced the ``atom of kinship'' in
his 1945 article ``Structural Analysis in Linguistics and in Anthropology.''
The atom comprises a man, his wife, their child, and the wife's brother: four
positions related by three relations (descent, marriage, avunculate).  Our
theorem captures the compositional closure that makes this atom stable.}

\begin{example}[Indian Caste Endogamy and Jati Compatibility]
The Indian caste system, particularly the institution of
\emph{jati} (birth-based sub-caste), provides one of the most elaborate
and long-lived matching systems in human history.  Each jati constitutes a
closed endogamous unit: marriage is permitted only within one's own jati,
or, in some reform movements, within a limited set of ``compatible'' jatis.

In our framework, each jati $j$ corresponds to an idea in $\IS$, encoding
the ritual status, occupational identity, dietary practices, and kinship
norms of the group.  Endogamy is the requirement that matching occurs only
between compatible ideas: $j_1 \res j_2$ holds only when $j_1$ and $j_2$
belong to marriageable categories.  Traditionally, the strictest form of
endogamy requires $j_1 = j_2$---self-matching---which is always guaranteed
by Theorem~\ref{thm:self-compat}.

The composition-preservation theorem (Theorem~\ref{thm:comp-preserves})
has a specific implication for the Indian system: if two jatis are
compatible and another pair of jatis is compatible, then the composed
cultural identities of their offspring are also compatible.  This ensures
the \emph{heritability} of caste compatibility across generations---precisely
the structural property that has allowed the jati system to persist for
millennia.  The system is self-reinforcing because compatibility, once
established, is algebraically guaranteed to reproduce itself through
composition.

B.\ R.\ Ambedkar's critique of caste endogamy---articulated in his 1916
paper ``Castes in India: Their Mechanism, Genesis and Development''---can be
read as a critique of this very algebraic closure.  Ambedkar argued that
caste is maintained by ``endogamy superimposed on exogamy'': the
prohibition of inter-caste marriage creates closed compatibility classes,
while the mandate of exogamy within the caste ensures that the compatibility
class reproduces itself.  Our formal framework captures exactly this
mechanism: self-compatibility guarantees the endogamous baseline, and
composition preservation guarantees hereditary closure.
\end{example}

\begin{theorem}[Right-Context Compatibility (12.6)]
\label{thm:right-context}
Compatible senders remain compatible when placed in the same right context:
\[
  a \res b
  \quad\Longrightarrow\quad
  (a \compose c) \res (b \compose c).
\]
\end{theorem}

\begin{proof}[Proof sketch]
By right-composition preservation of resonance.
\end{proof}

If two myths resonate, adding the same ritual context to both preserves their
resonance.  Shared ritual frames maintain cultural compatibility---which is
why liturgical uniformity (the Latin Mass, the Pali Canon) serves as a
\emph{resonance stabilizer} across geographically dispersed communities.

\begin{theorem}[Left-Context Compatibility (12.7)]
\label{thm:left-context}
Compatible ideas remain compatible when preceded by the same context:
\[
  a \res b
  \quad\Longrightarrow\quad
  (c \compose a) \res (c \compose b).
\]
\end{theorem}

\begin{proof}[Proof sketch]
By left-composition preservation of resonance.
\end{proof}

A shared cultural preamble---educational background, initiation rite---preserves
the resonance between ideas.  Preliminary socialization ensures continued
compatibility, which is why initiation rites (Masonic degrees, academic
comprehensive exams, military basic training) serve as \emph{compatibility
guarantors}: by imposing a common left-context $c$, they ensure that
subsequent ideas $a, b$ remain resonant in the initiates' minds.

\begin{example}[Academic Peer Review as Cultural Matching]
The peer review system in academic publishing instantiates a matching
problem with direct parallels to our formal framework.  When a journal
editor receives a manuscript (sender signal $s$), they must match it with
reviewers whose expertise (receiver repertoire) is compatible.  A paper
on Bayesian phylogenetics must be matched with a reviewer who can produce
a deep interpretation---$\depth(r_{\text{expert}} \compose s_{\text{paper}})$
is high---rather than a reviewer whose interpretation is shallow or null.

The left-context compatibility result (Theorem~\ref{thm:left-context})
applies directly: if two papers $a, b$ on related topics resonate
($a \res b$), then a reviewer with any background $c$ produces resonant
interpretations of both: $(c \compose a) \res (c \compose b)$.  This is
why reviewers are typically assigned clusters of related papers---the
resonance between papers ensures consistent evaluation.

The right-context result (Theorem~\ref{thm:right-context}) applies to
reviewer training: if two reviewers $c_1, c_2$ share a common disciplinary
background (are both ``preceded by'' the same graduate training $c$), then
their compatibility with any given paper is preserved.  The
standardization of doctoral training within a discipline serves as the
``shared left-context'' that ensures reviewer compatibility with the
discipline's literature.
\end{example}

% ----------------------------------------------------------------
\section{Quality Bounds}
\label{sec:ch12-quality}
% ----------------------------------------------------------------

\begin{theorem}[Composition Quality Bound (12.8)]
\label{thm:quality-bound}
The quality of any match is bounded by the sum of depths:
\[
  \operatorname{quality}(s, r) \;\le\; \depth(s) + \depth(r).
\]
\end{theorem}

\begin{proof}[Proof sketch]
Direct application of the subadditivity axiom.
\end{proof}

The depth of cultural synthesis---marriage, trade, dialogue---cannot exceed
the combined depth of the participants.  Two shallow cultures produce a
shallow synthesis.  This constrains utopian visions of ``fusion'': the
Esperanto of cultures is bounded above by the complexity of the cultures
being fused.

\begin{theorem}[Void Match Quality, Right (12.9)]
\label{thm:void-quality-right}
Matching with void yields one's own depth: $\operatorname{quality}(a, \void) = \depth(a)$.
\end{theorem}

\begin{proof}[Proof sketch]
$a \compose \void = a$ by the void axiom.
\end{proof}

The ``null exchange''---giving silence, receiving nothing---reveals one's own
depth.  This is the formal basis of introspection: in the absence of external
input, one sees only one's own complexity.  The contemplative
traditions---monastic silence, Zen \emph{zazen}, Quaker worship---formalize
this structural fact as spiritual practice.

\begin{theorem}[Void Match Quality, Left (12.10)]
\label{thm:void-quality-left}
The void sender reveals the receiver's depth: $\operatorname{quality}(\void, a) = \depth(a)$.
\end{theorem}

\begin{proof}[Proof sketch]
$\void \compose a = a$ by the void axiom.
\end{proof}

A culture that sends nothing (isolationism) reveals only what the receiver
already has.  The ``mirror effect'' of cultural isolation: the isolated
society, receiving no external signals, can only see its own reflection.

\begin{theorem}[Void-Void Match (12.11)]
\label{thm:void-void}
The quality of matching void with void is zero: $\operatorname{quality}(\void, \void) = 0$.
\end{theorem}

\begin{proof}[Proof sketch]
$\void \compose \void = \void$, and $\depth(\void) = 0$.
\end{proof}

When neither party brings anything, nothing is produced.  The ``empty
exchange'' has zero cultural value---the structural null point of all social
interaction.  This is Hobbes's state of nature, formalized: before the
social contract introduces non-void ideas, cultural quality is identically
zero.

% ----------------------------------------------------------------
\section{Match Pair Properties}
\label{sec:ch12-pairs}
% ----------------------------------------------------------------

\begin{theorem}[Compatible Pairs are Self-Evident (12.12)]
\label{thm:self-match}
The match pair $(a, a)$ is always compatible.
\end{theorem}

\begin{proof}[Proof sketch]
By self-compatibility (Theorem~\ref{thm:self-compat}).
\end{proof}

Self-matching (narcissistic culture) is always structurally valid, even if
socially prohibited.  The incest taboo in Lévi-Strauss's theory exists
precisely to prevent the structurally trivial but socially destructive
self-match.  Without the taboo, every lineage would choose itself---the
path of least structural resistance.

\begin{theorem}[Void Pair is Compatible (12.13)]
\label{thm:void-pair}
The match pair $(\void, \void)$ is compatible.
\end{theorem}

\begin{proof}[Proof sketch]
By self-compatibility of void.
\end{proof}

The ``null marriage''---two culturally empty entities exchanging nothing---is
trivially compatible.  This is the structural zero of kinship: before culture
introduces non-void ideas, all matches are vacuously valid.

\begin{theorem}[Composition Quality Associativity (12.14)]
\label{thm:quality-assoc}
Quality is associative under composition:
\[
  \operatorname{quality}((a \compose b), c)
  \;=\;
  \operatorname{quality}(a, (b \compose c)).
\]
\end{theorem}

\begin{proof}[Proof sketch]
Both sides unfold to $\depth((a \compose b) \compose c) = \depth(a \compose (b \compose c))$
by associativity.
\end{proof}

In a three-party exchange (mediator between two clans), the total quality is
invariant under regrouping.  The ``middle-man'' can be absorbed into either
side without changing the structural outcome---a formal justification for the
role of the \emph{go-between} in cross-cultural negotiation.

\begin{remark}[Open Questions on Dynamic Compatibility]
The matching theory developed in this chapter is fundamentally
\emph{static}: compatibility is a fixed relation on $\IS$, and match
quality is a fixed function of depth.  In practice, cultural compatibility
is dynamic---it evolves over time as matched pairs interact, adapt, and
transform each other.  Several open questions arise:

\begin{enumerate}
\item \textbf{Compatibility drift.}  Can two initially compatible ideas
  become incompatible through repeated composition?  Our framework precludes
  this (composition preserves compatibility by Theorem~\ref{thm:comp-preserves}),
  but this may reflect a limitation of the model rather than a feature of
  cultural reality.  Marriages do fail; alliances do dissolve.  A richer
  model might introduce a dynamic compatibility relation
  $\res_t$ that evolves with time.

\item \textbf{Matching under incomplete information.}  In the
  Gale-Shapley framework, agents know their complete preference orderings.
  In cultural matching, compatibility is often discovered only through
  interaction---one does not know whether two ideas resonate until one
  composes them.  A model incorporating costly compatibility testing
  (``courtship'') would connect to the economics of search and inspection
  goods.

\item \textbf{Many-to-many matching.}  Our match pairs are dyadic, but
  many cultural institutions involve polyadic matching: a student matched
  with multiple mentors, a polycultural society integrating many traditions
  simultaneously.  Extending the framework to $k$-tuples rather than pairs
  would require generalizing the compatibility and quality notions.

\item \textbf{Stability under perturbation.}  Is a stable matching robust
  to small changes in the compatibility relation?  If one pair's resonance
  is disrupted (e.g., by cultural change), does the entire matching unravel,
  or can it be locally repaired?  This connects to the literature on robust
  matching in mechanism design.
\end{enumerate}
\end{remark}

% ----------------------------------------------------------------
\section{Structural Matching Theorems}
\label{sec:ch12-structural}
% ----------------------------------------------------------------

\begin{theorem}[Double Composition Quality Bound (12.15)]
\label{thm:double-bound}
The quality of matching two compound ideas is quadruply bounded:
\[
  \operatorname{quality}((a \compose b), (c \compose d))
  \;\le\;
  \depth(a) + \depth(b) + \depth(c) + \depth(d).
\]
\end{theorem}

\begin{proof}[Proof sketch]
Apply subadditivity to the outer composition, then to each inner
composition, and combine via transitivity.
\end{proof}

The synthesis of two compound ideas---the alliance of two already-complex
cultural systems---cannot exceed the total depth of all components.  Colonial
encounters between complex civilizations (Aztec--Spanish, Chinese--British)
produce syntheses bounded by the sum of all four sub-cultural depths involved.

\begin{theorem}[Iterated Matching Depth (12.16)]
\label{thm:iterated-match}
Matching a signal with itself $n$ times via $\compN{a}{n}$ yields depth at
most $n \cdot \depth(a)$.
\end{theorem}

\begin{proof}[Proof sketch]
By the iterated composition depth bound (Theorem~\ref{thm:iterated-signal}
from Chapter~\ref{ch:marketplace}).
\end{proof}

Repeated cultural exchange---annual rituals, recurring trade---has bounded
cumulative depth.  The Kula ring's repeated exchanges are subadditive: each
cycle adds at most $\depth(\text{gift})$ to the cultural synthesis.
Malinowski's observation that the Kula values lie not in the objects
themselves but in the circulation is consistent with our bound: the \emph{same}
gift, circulated $n$ times, yields at most $n \cdot \depth(\text{gift})$,
not exponentially growing prestige.\footnote{Bronisław Malinowski,
\textit{Argonauts of the Western Pacific} (1922).}

\begin{theorem}[Synthesis Self-Compatibility (12.17)]
\label{thm:synthesis-self}
If $a$ and $b$ are compatible, then their composition $a \compose b$ is
compatible with itself.
\end{theorem}

\begin{proof}[Proof sketch]
Every idea is self-compatible (Theorem~\ref{thm:self-compat}).
\end{proof}

A culturally valid marriage (compatible pair) produces offspring (composition)
that is internally consistent.  The synthesis is never self-contradictory
in the resonance sense---a formal guarantee of the coherence of
cross-cultural offspring.

\begin{theorem}[Exchange Symmetry of Quality (12.18)]
\label{thm:exchange-symm}
If $a \res b$, then both $(a \compose b)$ and $(b \compose a)$ are
self-compatible:
\[
  a \res b
  \;\Longrightarrow\;
  (a \compose b) \res (a \compose b)
  \;\wedge\;
  (b \compose a) \res (b \compose a).
\]
\end{theorem}

\begin{proof}[Proof sketch]
Both directions follow from self-compatibility.
\end{proof}

In Mauss's \textit{Essai sur le don}, the gift and counter-gift are
structurally parallel but not identical.  We have $a \compose b \ne b \compose a$
in general---composition is not commutative---formalizing the asymmetry inherent
in reciprocal exchange.  The gift carries the giver's stamp; the counter-gift
carries the recipient's.  Yet both compositions are self-consistent: neither the
gift nor the counter-gift contradicts itself, even if they differ from each
other.  Reciprocity does not require identity---only mutual
self-consistency.\footnote{Marcel Mauss, \textit{Essai sur le don} (1925).
The \emph{hau}---the spirit of the gift that compels return---need not be
identical to the original gift, only resonant with it.}

\begin{lstlisting}[caption={Exchange symmetry in Lean~4}]
theorem exchange_compositions_self_compatible {a b : I}
    (h : compatible a b) :
    compatible (compose a b) (compose a b) /\
    compatible (compose b a) (compose b a) :=
  <self_compatible _, self_compatible _>
\end{lstlisting}

\begin{remark}[Connection to Chapter~\ref{ch:marketplace}: From Competition to Cooperation]
Chapters~\ref{ch:marketplace} and~\ref{ch:matching} are structural
complements.  In Chapter~\ref{ch:marketplace}, multiple senders compete
for a single receiver's interpretive engagement; the central question is
\emph{which signal dominates?}  In this chapter, senders and receivers are
matched pairwise; the central question is \emph{which pairings are stable?}

The formal connection is precise.  The dominance relation of
Chapter~\ref{ch:marketplace}---$s_1$ dominates $s_2$ on repertoire
$\mathbf{R}$ when
$\depth(r \compose s_1) \ge \depth(r \compose s_2)$ for all
$r \in \mathbf{R}$---can be interpreted as a preference ordering: the
receiver ``prefers'' $s_1$ to $s_2$.  When this preference ordering is
used to construct a Gale-Shapley instance, the resulting stable matching
is the \emph{equilibrium of the marketplace}.  Competition
(Chapter~\ref{ch:marketplace}) determines preferences; matching
(this chapter) determines allocation.

This connection illuminates a fundamental tension in cultural systems.  The
marketplace rewards \emph{breadth}---signals that achieve non-zero depth
across many receivers.  But matching rewards \emph{depth}---pairings that
achieve maximal composition quality for specific pairs.  A society that
optimizes for marketplace competition (broadcasting uniform signals) may
sacrifice matching quality (deep dyadic synthesis), and vice versa.  The
tension between broadcast media (marketplace logic) and intimate
relationships (matching logic) is a structural feature of ideatic space,
not merely a sociological observation.
\end{remark}

\subsection{Matching, Solidarity, and Cultural Reproduction}

The eighteen theorems of this chapter collectively establish that the
algebraic structure of ideatic space is \emph{hospitable to matching}: the
properties of reflexivity, symmetry, composition-preservation, and bounded
quality ensure that stable pairings exist, persist across generations, and
produce bounded but non-trivial cultural synthesis.  This hospitality is
not accidental.  The IDT axioms were chosen, in part, because they
guarantee exactly these matching-theoretic properties.  The axioms encode,
at the most abstract level, the structural preconditions for the
reciprocal exchange that Mauss identified as the foundation of social life.

Durkheim's distinction between \emph{mechanical solidarity} (solidarity
based on similarity) and \emph{organic solidarity} (solidarity based on
complementarity) maps onto two extreme regimes of our matching framework.
In mechanical solidarity, all ideas are mutually resonant
($a \res b$ for all $a, b \in \IS$), and every match is trivially
compatible.  The compatibility graph is complete, and any matching is stable.
In organic solidarity, compatibility is selective---$a \res b$ holds only
for specific pairs---and matching requires careful allocation.  The
transition from mechanical to organic solidarity, which Durkheim associated
with the division of labor, corresponds formally to the sparsification of
the compatibility graph: fewer edges, more constrained matching, more
complex equilibria.

The cultural reproduction of matching systems---their persistence across
generations---is guaranteed by composition-preservation
(Theorem~\ref{thm:comp-preserves}).  If this generation's matches are
compatible, then the ``offspring'' (compositions) of those matches are also
compatible with each other.  The matching system reproduces itself through
algebraic closure.  This is the formal mechanism underlying what Bourdieu
called \emph{reproduction}: the tendency of social structures to perpetuate
themselves through the very practices they engender.\footnote{Pierre
Bourdieu, \textit{La Reproduction} (1970), co-authored with Jean-Claude
Passeron.}

% ----------------------------------------------------------------
\section{Conclusion}
\label{sec:ch12-conclusion}
% ----------------------------------------------------------------

The matching theory developed in this chapter reveals that cultural
compatibility---defined as resonance---possesses exactly the algebraic
properties needed for stable pairing.  Compatibility is reflexive
(Theorem~\ref{thm:self-compat}), symmetric
(Theorem~\ref{thm:compat-symm}), and preserved under composition
(Theorem~\ref{thm:comp-preserves}).  These are not accidental features but
structural necessities: without them, the kinship systems catalogued by
Lévi-Strauss, the gift economies described by Mauss, and the patron--client
networks analyzed by Bourdieu would lack the algebraic foundation to persist
across generations.

Quality bounds (Theorem~\ref{thm:quality-bound}) ensure that no match
produces unbounded cultural synthesis, and void-match results
(Theorems~\ref{thm:void-quality-right}--\ref{thm:void-void}) establish the
baseline against which all exchange is measured.  In the next chapter, we
move from dyadic matching to \emph{coalitional} composition: the
cooperative game theory of cultural groups, formalized through Durkheim's
concept of solidarity.
